\subsection{Localized eigenstates} \label{ssection:quasimodes} We now construct specific eigenstates for the operator $k^{-2}\Delta_k$ with adapted localization properties.

\subsubsection{Energy $E=0$} We first investigate the low energy eigenstates in the semiclassical regime
\[
k^{-2}\Delta_k u_k = E_k u_k, \qquad u_k \in C^\infty(\Sigma,L^k),\; \|u_k\|_{L^2}=1,\; E_k = \mc{O}(1/k).
\]
This corresponds the regime of eigenvalues $\lambda_{k,m}$ with $k \to \infty$ and $m \leqslant C$, where $C > 0$ is a fixed constant. From now on, we~use the notation $h = 1/k$. Let $\underline{x} \in \Si$, $u \in L_{\underline{x}}$ with $|u|=1$ and define the Dirac section of
$L^k$ by $(\delta_u^k, s) = s(\underline{x})/u^k$ for any $k \geq 0$, $s \in C^\infty (\Sigma,L^k)$. We~introduce the section
\begin{gather} \label{eq:def_etat_coherent}
\bme_{k,m,\underline{x}} := \Pi_{k,m} \delta_u^k.
\end{gather}

The following holds:

\begin{proposition}
\label{proposition:e0}
Let $C > 0$ and $f \in C^\infty_{{\comp}}(T^* \Sigma)$. Then for all $k \geq 0$, $0 \leq m \leq C$, $\bme_{k,m,\underline{x}} \in C^\infty(\Sigma,L^k)$ satisfies
\begin{xalignat}{2}
\begin{split} \label{eq:estim1}
\| \bme_{k,m,\underline{x}} \|^2 &= \frac{k}{2\pi} (1 + \bigo (k^{-1/2})), \\
\langle \Op_k(f) \bme_{k,m,\underline{x}}, \bme_{k,m,\underline{x}} \rangle &= \| \bme_{k,m,\underline{x}} \|^2 \bigl(f (\underline{x},0) + \bigo
(k^{-1/2}) \bigr).
\end{split}
\end{xalignat}
\end{proposition}

As we will see, $\bme_{k,m,\underline{x}}$ is a Gaussian state centered at
$\underline{x}$. We~provide a proof building on \cite{Landau1,Landau2}.

\begin{proof}
The section $ \bme_{k,m,\underline{x}}$ is given in terms of the Schwartz kernel of
the projector $\Pi_{k,m}$ by
\begin{gather} \label{eq:noyau_coherent}
\bme_{k,m,\underline{x}} (x) \otimes \con{u}^k = \Pi_{m,k} (x, \underline{x})
\in L_x^k \otimes \con{L}_{\underline{x}}^k.
\end{gather}
It was proved in \cite{Landau1,Landau2} that $\Pi_{k,m} (x,\underline{x})$ has an
asymptotic expansion, which leads to the following expansion for $ \bme_{k,m,\underline{x}}$. Introduce coordinates $(x_1, x_2)$ centered at $\underline{x}$ and such that the metric $g$
is $(1 + \bigo (|x|)) (dx_1^2 + dx_2^2)$ with $|x|^2 = x_1^2 + x_2^2$. Choose a primitive $\al = \al_1
dx_1 + \al_2 dx_2 $ of the Riemannian volume $\vol$ on a neighborhood of $y$, such that $ \al_1 (\underline{x})
= \al_2 (\underline{x}) =0$ and $\partial_{x_i} \al_j (\underline{x})$ is antisymmetric. Let $t$ be a
frame of $L$ such that $t (\underline{x}) = u$ and $\nabla t = \tfrac{1}{i}
\al \otimes t$. Then one has on a neighborhood of~$\underline{x}$:
\begin{equation} \label{eq:theta}
\bme_{k,m,\underline{x}} (x) = \frac{k}{2 \pi} e^{-\frac{1}{4} k |x|^2 } \sum_{\ell
=0}^N k^{ -\sfrac{\ell}{2} } a_\ell (k^{\sfrac{1}{2} } x_1, k^{\sfrac{1}{2}}
x_2) t^k(x) + \bigo (k^{- \psfrac{N+1}{2}}),
\end{equation}
where the $\bigo $ is in the $\mathcal{C}^0$-norm and the $a_{\ell}$'s are polynomials depending on $m$ and smoothly on
$\underline{x}$. The leading coefficient $a_0$ is given in terms of the $m$-th
Laguerre polynomial $Q_m$ by
\begin{gather} \label{eq:leading}
a_0 (x_1, x_2) = Q_m (|x|^2), \qquad Q_m (t) = \frac{1}{m!} \Bigl(\frac{d}{dt}-1\Bigr)^m t^m.
\end{gather}
Moreover, $\bme_{k,m,\underline{x}}$ is $\bigo_{C^\infty} (k^{-\infty})$ on any compact set not
containing $\underline{x}$.
Formulas \eqref{eq:theta}, \eqref{eq:leading} correspond to \cite[Eqs.\,(26) \& (40)]{Landau1}.

Since $\Pi_{k,m}$
is a projector, $\| \bme_{k,m,\underline{x}} \|^2 = \Pi_{k,m}
(\underline{x},\underline{x})$, so~by \eqref{eq:theta},
\[
\| \bme_{k,m,\underline{x}} \|^2 = \frac{k}{2\pi} + \bigo (k^{1/2}),
\]
because $Q_m (0) =1$. We~then compute the asymptotic expansion of the scalar product
\[
\langle \Op_k(f) \bme_{k,m,\underline{x}},
\bme_{k,m,\underline{x}} \rangle = \sum_{\ell, \ell' \geqslant 0 }
k^{-\frac{1}{2} (\ell + \ell')} I_{\ell, \ell'}(f,k)
\]
with $I_{\ell, \ell'} (f,k)$ given by the
integral
\[
\Bigl(\frac{k}{2 \pi} \Bigr)^4 \!\!\int\! e^{ik (x- y) \xi - \frac{1}{4}k (|x|^2 + |y^2|) }
f (\tfrac{1}{2} (x+y), \xi) a_{\ell} (k^{\sfrac{1}{2}} x) \overline{
a_{\ell'} (k^{\sfrac{1}{2}} y)} \; \dd\xi_1 \dd\xi_2 \dd x_1 \dd x_2 \dd y_1 \dd y_2.
\]
Notice that the phase is quadratic, non degenerate. It follows from a stationary phase computation that $I_{\ell, \ell'} (f,k) = k C_{\ell, \ell'} f(0, 0) + \bigo
(k^{\sfrac{1}{2}})$ for some complex coefficient $C_{\ell, \ell'}$. This yields:
\[
\langle \Op_k(f) \bme_{k,m,\underline{x}}, \bme_{k,m,\underline{x}} \rangle
= k f(0,0) C_{0,0} + \bigo (k^{1/2})\]
and $C_{0,0} = (2 \pi)^{-1} $ because of the previous estimate of $\| \bme_{k,m,\underline{x}}
\|^2$. This concludes the proof. We~could actually slightly improve the result
by replacing the $\bigo (k^{-1/2})$'s in~\eqref{eq:estim1} by $\bigo (k^{-1})$'s, which follows
from the fact that
each $a_\ell$ has the same parity as~$\ell$.
\end{proof}

\subsubsection{Energies $0 < E < E_c$} We now investigate the eigenstates such that
\[
k^{-2}\Delta_{k} u_k = (E+o(1))u_k,
\]
with $0 < E < E_c$. Equivalently, this corresponds to the regime of
eigenvalues $\lambda_{k,m}$ with $k \to \infty$ and $\eps k B \leqslant m \leqslant (1- \eps)kB$, where $\eps> 0$ is fixed. {\color{black}Below, recall that $a$ is the principal symbol of $\bA_k$, see \eqref{equation:mercredi-aprem}.}

\begin{proposition} \label{lem:gaussian-beam}
Let $z_0 \in D^*_0\Sigma := D^*\Sigma\setminus \{ \xi = 0 \}$. Then, there exists a neighborhood $U \subset D^*_0\Sigma$ of $z_0$ such that for all $f \in C^\infty_{{\comp}}(T^*\Sigma)$, for all $k \geq 0$, $\eps kB \leq m \leq (1-\eps)kB$, and $z \in U$ such that
\[
m=ka(z),
\]
there exists $\bme_{k,m,z} \in C^\infty(\Sigma,L^{k})$ such that $\| \bme_{k,m,z} \|^2 = 1$ and the
projection
\[
\bmf_{k,m,z} := \Pi_{k,m} \bme_{k,m,z}
\]
satisfies
\begin{xalignat}{2}
\begin{split}
\| \bmf_{k,m,z} \|^2 &= \tfrac{1}{ 2 \sqrt \pi} k^{-\sfrac{1}{2}}(1+o_{f,U}(1)), \\
\langle \Op_k(f) \bmf_{k,m,z}, \bmf_{k,m,z} \rangle &= \|\bmf_{k,m,z}\|^2 \bigl(\langle f \rangle (z) + \bigo_{f,U}
(k^{-\sfrac{1}{4}}) \bigr).
\end{split}
\end{xalignat}
The remainder terms $\bigo$'s are uniform with respect to $k \geq 0$.
\end{proposition}

The state $\bme_{k,m,z}$ is a coherent state centered at $z$ and $\bmf_{k,m,z}$
is a Gaussian beam supported on the magnetic geodesic of $z$ (i.e. the flow line of $H_a^\Omega$ passing through $z$).

\begin{proof}
We use $h=1/k$ in the proof once again. Let $\bme_{k,m,z} \in C^\infty(\Sigma,L^k)$ be a Gaussian state centered at $z$ and such that $\|\bme_{k,m,z}\|_{L^2}=1$ (see \cite[Ex.\,1, Ch.\,5]{Zworski-12} for instance). A~quick computation using the stationary phase lemma reveals
\begin{gather} \label{eq:etat_coherent}
\Op_k(g) \bme_{k,m,z} = g (z) \bme_{k,m,z} + \bigo_{L^2} (
h^{\sfrac{1}{2}})
\end{gather}
for any $g \in C^\infty_{{\comp}}(T^*\Sigma)$. Then $\bmf_{k,m,z} := \Pi_{k,m} \bme_{k,m,z}$ (with $m=ka(z)$) satisfies by Lemma
\ref{lem:commutateur_espace_propre} that
\begin{xalignat}{2} \notag
\langle \Op_k(f) \bmf_{k,m,z}, \bmf_{k,m,z} \rangle & = \langle \Op_k(\langle f \rangle)
\bme_{k,m,z}, \Pi_{k,m} \bme_{k,m,z} \rangle + \bigo (h \|\bmf_{k,m,z}\|) \\
& = \langle f \rangle (z) \|\bmf_{k,m,z}\|^2 + \bigo (h^{\sfrac{1}{2}} \|\bmf_{k,m,z}\|).
\end{xalignat}
To estimate the norm of $\bmf_{k,m,z}$, we~use
\eqref{eq:integrale_projecteur}, namely (recall that $U_t = e^{itk\bA_k}$):
\begin{gather} \label{eq:int_norme}
\begin{split}
\|\bmf_{k,m,z}\|^2 = \langle \Pi_{k,m} \bme_{k,m,z}, \bme_{k,m,z} \rangle & = \frac{1}{2 \pi }
\int_0^{2 \pi } e^{-itm } \langle U_t \bme_{k,m,z},\bme_{k,m,z}
\rangle \; \dd t \\
& = \frac{1}{2 \pi } \int_0^{2 \pi } \langle e^{itk(\bA_k-a(z))}
\bme_{k,m,z},\bme_{k,m,z} \rangle \; \dd t,
\end{split}
\end{gather}
where we used that $m = k a(z)$.
By assumption, the microsupport of $\bme_{k,m,z}$ is $\{z\}$ (that is for
any $\varphi \in C^\infty _{{\comp}} (T^* \Si)$ identically equal to 1 on a
neighborhood of $z$, $\Op_k(\varphi) \bme_{k,m,z} = \bme_{k,m,z} + \bigo_{C^\infty} (
h^{\infty})$). Then by Egorov's theorem, the microsupport of
$U_t \bme_{k,m,z}$ is $ \{ \Phi_t^a(z) \}$. Since $\Phi_t^a (z) = z $ if and only if $t$
is a multiple of $2 \pi$, the integral
\eqref{eq:int_norme} is modified by a $\bigo (h^{\infty})$ if we
restrict the integral to a neighborhood of $t=0$.

Microlocally, $\bA_k-a(z)$ can be conjugated by a
Fourier Integral Operator to $-ih\partial_{x_1}$ acting on
$\R^2_{x_1, x_2}$, see \cite[Th.\,12.3]{Zworski-12}. Since we can restrict the integral \eqref{eq:int_norme} to a
neighborhood of the origin, it~is sufficient to do the computation in this
model. The propagator is then given by $e^{itk(\bA_k-a(z))} \Psi (x_1, x_2) = \Psi (x_1 +t, x_2)$ for $\Psi \in C^\infty(\R^2)$. The point $z$ is in the characteristic of $\bA_k-a(z)$; it is therefore mapped to \hbox{$(\underline{x}, \underline{\xi}_1=0,\underline{\xi}_2) \in T^*\R^2$}. Up to conjugating again by a translation, we~can further assume that $\underline{x}=0$. Let
$\Psi_{\underline{x}, \underline{\xi}} \in C^\infty (\R^2)$ be
the coherent state
\[
\Psi_{\underline{x}, \underline{\xi}} (x) = (\pi h)^{-1/2} e^{ -
\frac{1}{2h} | x - \underline{x}|^2 + \frac{i}{h} \underline{\xi} (x - \underline{x}) } = (\pi h)^{1/2}e^{ -
\frac{1}{2h} | x|^2 + \frac{i}{h} \underline{\xi}_2 x_2 }.\]
By an exact computation, we~find that for $\delta > 0$, $\chi \in C^\infty_{{\comp}}(-\delta,\delta)$ such that $\chi \equiv 1$ in a neighborhood of $0$,
\[
\begin{split} \frac{1}{2 \pi } \int_0^{2 \pi } \langle e^{itk(\bA_k-a(z))} \bme_{k,m,z}&,\bme_{k,m,z} \rangle_{L^2} \chi(t)  \dd t \\
& = \frac{1}{2\pi} \int_{-\infty}^{+\infty} \int_{\R^2} \Psi _{ \underline{x},
\underline{\xi} } (x_1 + t, x_2) \con { \Psi_{ \underline{x}, \underline{\xi} } (x_1,
x_2)} \chi(t) \; \dd x_1 \dd x_2 \dd t \\
& = \frac{ h^{1/2}}{2 \sqrt \pi }\,(1+ \mc{O}(h^\infty)).
\end{split}
\]
Going back to \eqref{eq:int_norme}, we~find that $ \|\bmf_{k,m,z}\|^2_{L^2} =
\tfrac{1}{ 2 \sqrt \pi}\, h^{1/2} + \bigo (h^{\infty})$, which
concludes the proof.
\end{proof}

\section{Proof of the main results}

\label{section:proofs}

In this final section, we~prove Theorems \ref{theorem:main}, \ref{theorem:constant2} and conclude with the proof of Theorem \ref{theorem:horocyclic}. Up to a preliminary rescaling, we~can always assume that $B=1/|\chi(\Sigma)|$.

\subsection{Proof of Theorem \ref{theorem:main}} We can now characterize the semiclassical defect measures in the three regimes.

\skpt
\begin{proof}[Proof of Theorem \ref{theorem:main}]

\eqref{theorem:maini}
Suppose $E=0$. We~claim that any probability measure $\mu$ on $\Sigma$ is a semiclassical limit of eigenstates \eqref{equation:eigenstates}. First, given $x \in \Sigma$,
\[
\bme'_{k,m,x} := \bme_{k,m,x}/\|\bme_{k,m,x}\|_{L^2}
\]
converges semiclassically to $\delta_x$ by Proposition \ref{proposition:e0}. Given $x_1,\dots, x_N \in \Sigma$, $\bE_{k,m,N} := N^{-1/2}(\bme'_{k,m,x_1} +\cdots+ \bme'_{k,m,x_N})$ converges semiclassically to $N^{-1}(\delta_{x_1}+\cdots+\delta_{x_N})$ and $\|\bE_{k,m,N}\|^2_{L^2} = 1 + \mc{O}(k^{-\infty})$. Indeed, to compute the norm, let $\varphi_{i} \in C^\infty(\Sigma)$ be a small bump function centered at $x_i$, equal to $1$ on a small neighborhood of $x_i$ and such that $x_j \notin \supp(\varphi_i)$ for $j \neq i$. It is then immediate to verify that $\varphi_i \bme'_{k,m,x_i} = \bme'_{k,m,x_i} + \mc{O}_{C^\infty}(h^\infty)$ and $\varphi_j \bme'_{k,m,x_i} = \mc{O}_{C^\infty}(h^\infty)$. This leads to
\[
\begin{split}
\|\bE_{k,m,N}\|^2_{L^2} & = N^{-1}\langle \bme'_{k,m,x_1} +\cdots+ \bme'_{k,m,x_N}, \bme'_{k,m,x_1} +\cdots+ \bme'_{k,m,x_N}\rangle_{L^2} \\
& = 1 + N^{-1} \sum_{i \neq j} \langle \bme'_{k,m,x_i},\bme'_{k,m,x_j}\rangle_{L^2} \\
& = 1 + N^{-1} \sum_{i \neq j} \langle \varphi_i \bme'_{k,m,x_i},\bme'_{k,m,x_j}\rangle_{L^2} + \mc{O}(h^\infty) \\
& = 1 + N^{-1} \sum_{i \neq j} \langle \bme'_{k,m,x_i},\varphi_i \bme'_{k,m,x_j}\rangle_{L^2} + \mc{O}(h^\infty) = 1 + \mc{O}(h^\infty).
\end{split}
\]
The same computation shows that the associated defect measure is equal to \hbox{$N^{-1}(\delta_{x_1}+\cdots+\delta_{x_N})$}. Finally, since Dirac masses are dense in probability measures, a diagonal extraction allows to construct a family $(\bE_{k_n,m_n,N_n})_{n \geq 0}$ with microlocal limit $\mu$, and that for any probability measure $\mu$ on $\Sigma$.

We now assume that $0 < E < E_c$. Fix $z \in \{p=E\}$ and let $(z_k)_{k \geq 0}$ be a sequence converging to $z$ such that $a(z_k) = k^{-1}m_k$ with $m_k \in \Z_{\geq 0}$. By Proposition~\ref{lem:gaussian-beam}, $\bmf_{k,m_k,z_k}/\|\bmf_{k,m_k,z_k}\|_{L^2}$ converges microlocally to $\delta_z$, the Dirac masses carried by the periodic orbit of $z$. As above, one can also realize any convex linear combination of such Dirac masses. As these are dense in flow-invariant probability measures on $\{p=E\}$, this proves the claim.

\smallskip

\eqref{theorem:mainii}
This follows immediately from Proposition \ref{proposition:scl-measure}, item \eqref{proposition:scl-measureii}, and the unique ergodicity of the horocyclic flow (see Theorem \ref{lemma:horocycle}, item \eqref{lemma:horocycleii}).

\smallskip

\eqref{theorem:mainiii}
The proof follows \emph{verbatim} the standard proof of Quantum Ergodicity. In~the framework of twisted quantization, this question was recently addressed in \cite[\S 5.3]{Cekic-Lefeuvre-24} (in a more general framework).
\end{proof}
